% ===== PRÉAMBULE CANONIQUE — à coller VERBATIM en tête de CHAQUE .tex =====
\documentclass[11pt,a4paper]{article}
\usepackage[utf8]{inputenc}\usepackage[T1]{fontenc}\usepackage{lmodern}
\usepackage[french]{babel}
\usepackage{amsmath,amssymb,mathtools}\usepackage{icomma}
\usepackage{siunitx}\sisetup{locale=FR}  % \qty{}{}, \si{}, \ang{}
\usepackage{xcolor}\usepackage{colortbl}
\usepackage{tikz}\usepackage{pgfplots}\pgfplotsset{compat=1.18}
\usepackage[siunitx,europeanresistors]{circuitikz} % circuits (élec.)
\usepackage[most]{tcolorbox}\usepackage{enumitem}\usepackage{array,booktabs}
\usepackage[a4paper,margin=2.2cm]{geometry}
\usetikzlibrary{arrows.meta,calc,angles,quotes,positioning,patterns,%
  backgrounds,shapes.geometric,decorations.pathmorphing,decorations.markings,intersections}
\definecolor{primary}{RGB}{222,49,99}\definecolor{primarydark}{RGB}{150,22,55}
\definecolor{defcol}{RGB}{0,114,178}\definecolor{methcol}{RGB}{52,92,148}
\definecolor{excol}{RGB}{0,135,90}\definecolor{attncol}{RGB}{200,40,55}
\tcbset{boxbase/.style={enhanced,breakable,boxrule=0.9pt,arc=2.5pt,
  left=3mm,right=3mm,top=2.3mm,bottom=2.3mm,fonttitle=\bfseries,coltitle=white}}
% --- boîtes TUTORIEL (noms EXACTS, ne pas renommer) ---
\newtcolorbox{cle}[1][]{boxbase,colback=defcol!6,colframe=defcol,
  title={\textsc{À retenir}\ifx\relax#1\relax\relax\else\textmd{ : \itshape #1}\fi}}
\newtcolorbox{methode}[1][]{boxbase,colback=methcol!7,colframe=methcol,sharp corners,
  title={\textsc{Méthode}\ifx\relax#1\relax\relax\else\textmd{ --- \itshape #1}\fi}}
\newtcolorbox{exemplebox}[1][]{boxbase,colback=excol!5,colframe=excol,
  title={\textsc{Exemple}\ifx\relax#1\relax\relax\else\textmd{ : \itshape #1}\fi}}
\newtcolorbox{piege}[1][]{boxbase,colback=attncol!6,colframe=attncol,
  title={\textsc{Piège}\ifx\relax#1\relax\relax\else\textmd{ : \itshape #1}\fi}}
\newtcolorbox{remarque}[1][]{boxbase,colback=black!4,colframe=black!45,
  title={\textsc{Remarque}\ifx\relax#1\relax\relax\else\textmd{ : \itshape #1}\fi}}
% --- boîtes EXERCICE / CORRIGÉ (noms EXACTS, auto-comptées) ---
\newtcolorbox[auto counter]{exo}[1][]{boxbase,colback=primary!4,colframe=primary,
  title={\textsc{Exercice}~\thetcbcounter\ifx\relax#1\relax\relax\else\textmd{ --- \itshape #1}\fi}}
\newtcolorbox[auto counter]{corrige}[1][]{boxbase,colback=excol!4,colframe=excol,
  title={\textsc{Corrigé}~\thetcbcounter\ifx\relax#1\relax\relax\else\textmd{ --- \itshape #1}\fi}}
\newcommand{\vv}[1]{\vec{#1}}\newcommand{\ds}{\displaystyle}
\newcommand{\R}{\mathbb{R}}\newcommand{\N}{\mathbb{N}}\newcommand{\Z}{\mathbb{Z}}
% =========================================================================
\begin{document}
\usetikzlibrary{babel}
\begin{corrige}[loi de Pouillet directe]
\[ R=\rho\,\frac{L}{S}=1{,}7\times10^{-8}\times\frac{50}{10^{-6}}=\qty{0.85}{\ohm}. \]
\end{corrige}

\begin{corrige}[penser à convertir la section]
\begin{enumerate}[label=\alph*)]
  \item $S=\qty{2.5}{\milli\metre\squared}=2{,}5\times10^{-6}\,\si{\metre\squared}$.
  \item $R=\rho\dfrac{L}{S}=2{,}8\times10^{-8}\times\dfrac{200}{2{,}5\times10^{-6}}=\qty{2.24}{\ohm}$.
\end{enumerate}
\end{corrige}

\begin{corrige}[proportionnalités]
\begin{enumerate}[label=\alph*)]
  \item $R\propto L$ : doubler $L$ \emph{double} $R$, soit $R=\qty{12}{\ohm}$.
  \item $R\propto\dfrac{1}{S}$ : tripler $S$ \emph{divise} $R$ par $3$, soit $R=\qty{2}{\ohm}$.
\end{enumerate}
\end{corrige}

\begin{corrige}[conductivité]
\[ \sigma=\frac{1}{\rho}=\frac{1}{1{,}7\times10^{-8}}\approx\qty{5.9e7}{\siemens\per\metre}. \]
L'unité est le siemens par mètre (\si{\siemens\per\metre}), soit \si{\per\ohm\per\metre}.
\end{corrige}

\begin{corrige}[résistance et température]
\[ R=R_0(1+\alpha T)=20\times\bigl(1+4\times10^{-3}\times100\bigr)=20\times(1+0{,}4)=\qty{28}{\ohm}. \]
\end{corrige}

\end{document}
